\documentclass[fontsize=12pt,paper=a4,twoside]{scrartcl}

\newcommand{\grad}{\ensuremath{^{\circ}} }
\renewcommand{\strut}{\vrule width 0pt height5mm depth2mm}

\usepackage[utf8]{inputenc}
\usepackage[final]{pdfpages}
% obere Seitenränder gestalten können
\usepackage{fancyhdr}
\usepackage{moreverb}
% Graphiken als jpg, png etc. einbinden können
\usepackage{graphicx}
\usepackage{stmaryrd}
% Floats Objekte mit [H] festsetzen
\usepackage{float}
% setzt URL's schön mit \url{http://bla.laber.com/~mypage}
\usepackage{url}
% Externe PDF's einbinden können
\usepackage{pdflscape}
% Verweise innerhalb des Dokuments schick mit " ... auf Seite ... "
% automatisch versehen. Dazu \vref{labelname} benutzen
\usepackage[ngerman]{varioref}
\usepackage[ngerman]{babel}
\usepackage{ngerman}
% Bibliographie
\usepackage{bibgerm}
% Tabellen
\usepackage{tabularx}
\usepackage{supertabular}
\usepackage[colorlinks=true, pdfstartview=FitV, linkcolor=blue,
            citecolor=blue, urlcolor=blue, hyperfigures=true,
            pdftex=true]{hyperref}
\usepackage{bookmark}
\usepackage{enumitem}

% Damit Latex nicht zu lange Zeilen produziert:
\sloppy
%Uneinheitlicher unterer Seitenrand:
%\raggedbottom

% Kein Erstzeileneinzug beim Absatzanfang
% Sieht aber nur gut aus, wenn man zwischen Absätzen viel Platz einbaut
\setlength{\parindent}{0ex}

% Abstand zwischen zwei Absätzen
\setlength{\parskip}{1ex}

% Seitenränder für Korrekturen verändern
\addtolength{\evensidemargin}{-1cm}
\addtolength{\oddsidemargin}{1cm}

\bibliographystyle{gerapali}

% Lustige Header auf den Seiten
  \pagestyle{fancy}
  \setlength{\headheight}{70.55003pt}
  \fancyhead{}
  \fancyhead[LO,RE]{Software-Projekt 2\\ WiSe 
  \\Testprotokoll}
  \fancyhead[LE,RO]{Seite \thepage\\\slshape \leftmark\\\slshape \rightmark}

%Etwas größere Zeilenabstände in Tabellen
\renewcommand{\arraystretch}{1.2}

%
% Und jetzt geht das Dokument los....
%

\begin{document}

% Lustige Header nur auf dieser Seite
  \thispagestyle{fancy}
  \fancyhead[LO,RE]{ }
  \fancyhead[LE,RO]{Universität Bremen\\FB 3 -- Informatik\\
  Prof. Dr. Rainer Koschke \\TutorIn: Karsten Hölscher}
  \fancyfoot[C]{}

% Start Titelseite
  \vspace{3cm}

  \begin{minipage}[H]{\textwidth}
  \begin{center}
  \bf
  \Large
  Software--Projekt 2 2014\\
  \smallskip
  \small
  VAK 03-BA-901.02\\
  \vspace{3cm}
  \end{center}
  \end{minipage}
  \begin{minipage}[H]{\textwidth}
  \begin{center}
  \vspace{1cm}
  \bf
  \Large Testprotokoll\\
  \vspace{3ex}
   	  \begin{figure}[H]
      \centering
      \includegraphics[width=0.25\textwidth]{../WOYM.png}
      \end{figure}
  \vfill
  \end{center}
  \end{minipage}
  \vfill
  \begin{minipage}[H]{\textwidth}
  \begin{center}
  \sf
  \begin{tabular}{l}
  Tim Hansen \\
  Adrian Lück \\
  Jurij Schmidt\\
  \end{tabular}

  \end{center}
  \end{minipage}

% Ende Titelseite

% Start Leerseite

\newpage

  \thispagestyle{fancy}
  \fancyhead{}
  \fancyhead[LO,RE]{Software-Projekt 2 \\  WiSe 2014
  \\Testprotokoll}
    \fancyhead[LE,RO]{Seite \thepage\\\slshape \mbox{}\\\slshape}
  \fancyfoot{}
  \renewcommand{\headrulewidth}{0.4pt}

\begin{tabular}{ll}
\textbf{Tester:} & Adrian Lück\\
\textbf{Datum:} & 07.02.2015\\
\textbf{Testfälle:} & 362 
\end{tabular}

Es wurden alle Unit- und Integrationstests ausgeführt, die in den jeweiligen XML-Dateien (\textit{unitTests.xml} bzw. \textit{integrationTests.xml}) aufgelistet sind. Alle Testfälle verliefen fehlerfrei. Es wurde eine Zweigabdeckung von 44,9 \% mithilfe des Eclipse-Plugins \textit{EclEmma} ermittelt.\\

Folgende Zweigabdeckungen wurden für die einzelnen Pakete ermittelt:

\renewcommand{\arraystretch}{1.3}
\begin{tabularx}{\textwidth}{|X|p{2cm}|X|}
\hline
\textbf{Paket} & \textbf{Zweig-abdeckung} & \textbf{Kommentar} \\\hline
\texttt{org.woym.common.config} & 88,4 \% & \\\hline
\texttt{org.woym.common.exceptions} & 100,0 \% & \\\hline
\texttt{org.woym.common.messages} & 81,0 \% & \\\hline
\texttt{org.woym.common.objects} & 73,8 \% &  \\\hline
\texttt{org.woym.controller} & 0,0 \% & Lediglich Aufrufe von bereits getesteten Methoden, wobei noch ggf. Exceptions gefangen werden. Daher keine Testfälle.\\\hline
\texttt{org.woym.controller.manage} & 0,0 \% & Viele getter- und setter-Methoden. Nicht vollständig triviale, aber doch einfache Operationen. Daher keine Testfälle. \\\hline
\texttt{org.woym.controller.planning} & 17,4 \% & Nur Abdeckung von umfangreicheren Methoden der Klasse \texttt{PlanningController}. Abdeckung von 36,4 \% für diese Klasse.\\\hline
\texttt{org.woym.logic} & 25,3 \% & Die geringe Abdeckung resultiert im Wesentlichen daraus, dass keine Testfälle für die Klasse \texttt{BackupRestoreHandler} existieren, da diese schwer umzusetzen wären. Der in dem Paket befindliche \texttt{CommandHandler} ist jedoch zu 95,3 \% abgedeckt.\\\hline
\texttt{org.woym.logic.command} & 100,0 \% & \\\hline
\end{tabularx}

\begin{tabularx}{\textwidth}{|X|p{2cm}|X|}
\hline
\texttt{org.woym.logic.spec} & 100,0 \% & \\\hline
\texttt{org.woym.logic.util} & 58,5 \% & Die sehr wichtige enthaltene Klasse \texttt{ActivityValidator} ist zu 100,0 \% abgedeckt.\\\hline
\texttt{org.woym.persistence} & 47,4 \% & Die Funktionalität aller Datenbankanfragen wird getestet. Die niedrige Testabdeckung entsteht, da null Prüfungen und das Fangen von Exceptions nicht getestet wurden. \\\hline
\texttt{org.woym.ui.converters} & 52,7 \% & \\\hline
\texttt{org.woym.ui.util} & 34,7 \% & Die wichtige Klasse \texttt{ActivityParser} ist zu 100,0 \% abgedeckt.\\\hline
\texttt{org.woym.ui.validators} & 20,1 \% & \\\hline
\end{tabularx}

\end{document} 
